\section*{Capítulo \ref{ch:g5:decimals} --- Los números decimales}

\begin{solution}{exo:g5:decimals:1}
$5.207$; \quad $0.85$; \quad $4.507$; \quad $12.009$.
\end{solution}

\begin{solution}{exo:g5:decimals:2}
El $4$ cuenta las décimas y el $8$, las milésimas:
\[
27.418 = 27 + \frac{4}{10} + \frac{1}{100} + \frac{8}{1000}.
\]
\end{solution}

\begin{solution}{exo:g5:decimals:3}
$3.60 = 3.6$: verdadero (un \omterm{def:g1:counting:zero}{cero} final no añade nada).

$0.5 = 0.05$: falso --- $5$ décimas frente a $5$ centésimas.

$2.070 = 2.07$: verdadero.

$1.3 = 1.03$: falso --- $3$ décimas frente a $3$ centésimas.
\end{solution}

\begin{solution}{exo:g5:decimals:4}
$4.8 > 4.79$; \quad $0.302 < 0.32$ (décimas iguales y centésimas
$0 < 2$); \quad $6.5 = 6.500$; \quad $9.09 < 9.1$.
\end{solution}

\begin{solution}{exo:g5:decimals:5}
$3.041 < 3.104 < 3.14 < 3.4 < 3.41$.
\end{solution}

\begin{solution}{exo:g5:decimals:6}
$6.83$: entre $6$ y $7$; entre $6.8$ y $6.9$.

$0.492$: entre $0$ y $1$; entre $0.4$ y $0.5$.

$12.06$: entre $12$ y $13$; entre $12.0$ y $12.1$.
\end{solution}

\begin{solution}{exo:g5:decimals:7}
$8.276$: a la unidad, $8$; a la décima, $8.3$; a la centésima,
$8.28$. Y $3.95$ a la décima: la cifra de las centésimas es $5$, así que redondea
hacia arriba: $4.0$.
\end{solution}

\begin{solution}{exo:g5:decimals:8}
Por ejemplo, $2.65$; \quad $0.415$; \quad $5.995$. (Hay infinitas
respuestas en cada caso.)
\end{solution}

\begin{solution}{exo:g5:decimals:9}
$0.85 < 0.9 < 1.02 < 1.2$ (en metros).
\end{solution}

\begin{solution}{exo:g5:decimals:10}
Hay que comparar primero las décimas: $7.12$ tiene $1$ décima y $7.9$ tiene $9$
décimas, así que $7.12 < 7.9$. Lucía comparó las partes decimales como números
enteros ($12$ frente a $9$), olvidando que $12$ centésimas es mucho menos
que $90$ centésimas.
\end{solution}

\begin{solution}{exo:g5:decimals:11}
Al redondear a la décima se obtiene $6.4$ para los números de dos decimales desde
$6.35$ (la \omterm{def:g3:division:half}{mitad} justa sube) hasta $6.44$: los números $6.35$, $6.36$,
\dots, $6.44$. El menor: $6.35$; el mayor: $6.44$.
\end{solution}
